\documentclass{subfiles}
\begin{document}
    Let \(m\) be the document \(\A\) produced, 
        of which we want to garantee the autenticity.
        Since signing the whole document is tedious, and most of the time useless,
        we can simply sign part of it, and to do so we consider the hash \(h(m)\) of the document.
        Here, \(h\) is what is called an hash function.

    Although the literature provides a pletora of things to say about hash functions,
        what matters here is simply understanding that an hash functions 
        maps elements from a certain universe \(U\) (e.g. the set of all possible messages),
        to the set \(\set{0, \ldots, m}\), for some integer \(m\).
        The important thing is that the number of collisions, i.e.,
        the number of elements of \(U\) such that the hash coicides,
        is the lowest possible. Again, the literature provide us with several of such functions,
        thus, here we can assume \(h\) one of such hash functions.

    Back to the algorithm; besides computing the hash of \(m\),
        the RSA algorithm proceed as described in \Cref{sec:RSA},
        with the difference being that instead of cyphering \(m\), 
        we cypher \(h(w)\).
\end{document}
